\documentclass[11pt,a4paper]{article}
\usepackage[utf8]{inputenc}
\usepackage[margin=0.8in]{geometry}
\usepackage{listings}
\usepackage{xcolor}
\usepackage{titlesec}
\usepackage{hyperref}
\usepackage{fancyhdr}

\definecolor{codebg}{RGB}{245,245,245}
\definecolor{codegreen}{RGB}{0,128,0}
\definecolor{codegray}{RGB}{128,128,128}
\definecolor{codeblue}{RGB}{0,0,180}
\definecolor{codepurple}{RGB}{128,0,128}

\lstset{
    language=C++,
    backgroundcolor=\color{codebg},
    basicstyle=\ttfamily\footnotesize,
    keywordstyle=\color{codeblue}\bfseries,
    commentstyle=\color{codegreen}\itshape,
    stringstyle=\color{codepurple},
    numberstyle=\tiny\color{codegray},
    numbers=left,
    numbersep=5pt,
    frame=single,
    breaklines=true,
    breakatwhitespace=false,
    tabsize=4,
    showstringspaces=false,
    captionpos=b,
    morekeywords={override,nullptr,string,shared_ptr,unique_ptr,make_shared,make_unique,
                  lock_guard,unique_lock,mutex,thread,vector,unordered_map,map,
                  chrono,steady_clock,time_point,atomic,condition_variable,
                  function},
}

\pagestyle{fancy}
\fancyhf{}
\fancyhead[L]{\textbf{LLD --- Job Scheduler}}
\fancyhead[R]{\thepage}

\title{\Huge\textbf{Job Scheduler}\\[0.3em]\Large Low-Level Design --- C++ Concurrency}
\author{LLD Interview Preparation}
\date{}

\begin{document}
\maketitle
\tableofcontents
\newpage

% ============================================================
\section{File Structure}
% ============================================================
\begin{verbatim}
job-scheduler/
  models.hpp       -- JobStatus, JobPriority, RetryPolicy,
                      Job (jobMtx, serialize/deserialize)
  strategies.hpp   -- SchedulingStrategy, JobObserver
  managers.hpp     -- JobScheduler (Singleton, worker pool,
                      persistence, dead letter queue)
  main.cpp         -- Demo: priorities, retries, backoff,
                      dead letter, persistence
\end{verbatim}

\textbf{Compilation:}
\begin{verbatim}
g++ -std=c++17 -pthread -o job_scheduler.exe main.cpp
\end{verbatim}

\subsection{Concurrency Summary}
\begin{itemize}
    \item \textbf{Level 1:} \texttt{schedulerMtx} (JobScheduler) --- pendingJobs, allJobs, deadLetterQueue, observers
    \item \textbf{Level 2:} \texttt{jobMtx} (per-Job) --- status, retry count, last error
    \item \textbf{condition\_variable:} workers block when no jobs; \texttt{submitJob} signals one worker
    \item \textbf{unique\_lock:} required for \texttt{cv.wait()} in worker loop
    \item \textbf{Exponential backoff:} delay = baseDelay $\times$ $2^{\text{attempt}}$, then re-queue
    \item \textbf{Dead letter queue:} after maxRetries exhausted
    \item \textbf{Persistence:} \texttt{serialize/deserialize} for survive-restart
\end{itemize}

\subsection{Design Patterns}
\begin{itemize}
    \item \textbf{Singleton:} JobScheduler
    \item \textbf{Strategy:} PriorityFirstStrategy, FIFOStrategy
    \item \textbf{Observer:} LoggingObserver, AlertObserver
\end{itemize}

\subsection{Smart Pointers}
\begin{itemize}
    \item \texttt{shared\_ptr<Job>} --- scheduler + workers + caller all hold refs
    \item \texttt{unique\_ptr<SchedulingStrategy>} --- scheduler sole owner
    \item \texttt{shared\_ptr<JobObserver>} --- scheduler vector + caller
\end{itemize}

% ============================================================
\newpage
\section{models.hpp}
% ============================================================
\begin{lstlisting}[title={\texttt{models.hpp}}]
#ifndef MODELS_HPP
#define MODELS_HPP

#include <string>
#include <vector>
#include <iostream>
#include <mutex>
#include <memory>
#include <chrono>
#include <functional>
#include <fstream>
#include <sstream>
#include <atomic>
using namespace std;

// ============================================================
// ENUMS
// ============================================================
enum class JobStatus {
    PENDING,
    RUNNING,
    COMPLETED,
    FAILED,
    DEAD_LETTERED
};

enum class JobPriority {
    LOW = 0,
    MEDIUM = 1,
    HIGH = 2,
    CRITICAL = 3
};

inline string statusToString(JobStatus s) {
    switch (s) {
        case JobStatus::PENDING:       return "PENDING";
        case JobStatus::RUNNING:       return "RUNNING";
        case JobStatus::COMPLETED:     return "COMPLETED";
        case JobStatus::FAILED:        return "FAILED";
        case JobStatus::DEAD_LETTERED: return "DEAD_LETTERED";
        default: return "UNKNOWN";
    }
}

inline string priorityToString(JobPriority p) {
    switch (p) {
        case JobPriority::LOW:      return "LOW";
        case JobPriority::MEDIUM:   return "MEDIUM";
        case JobPriority::HIGH:     return "HIGH";
        case JobPriority::CRITICAL: return "CRITICAL";
        default: return "UNKNOWN";
    }
}

// ============================================================
// RETRY POLICY
// Defines max retries and base delay for exponential backoff.
// Immutable after creation -- no mutex needed.
// ============================================================
class RetryPolicy {
    int maxRetries;
    int baseDelayMs;   // base delay in milliseconds

public:
    RetryPolicy(int maxRetries = 3, int baseDelayMs = 100)
        : maxRetries(maxRetries), baseDelayMs(baseDelayMs) {}

    int getMaxRetries() const { return maxRetries; }
    int getBaseDelayMs() const { return baseDelayMs; }

    // Exponential backoff: delay = base * 2^attempt
    int getDelayMs(int attempt) const {
        return baseDelayMs * (1 << attempt);
    }
};

// ============================================================
// JOB
// Has its own mutex for thread-safe status updates.
// Workers lock the job while executing it.
// PESSIMISTIC: lock job before any state change.
//
// Persist format: "id|name|priority|status|retryCount|maxRetries|baseDelay"
// ============================================================
class Job {
    string jobId;
    string name;
    JobPriority priority;
    JobStatus status;
    int retryCount;
    RetryPolicy retryPolicy;
    function<bool()> task;  // the actual work to execute
    mutable mutex jobMtx;
    chrono::steady_clock::time_point scheduledTime;
    string lastError;

public:
    Job(const string& id, const string& name, JobPriority priority,
        RetryPolicy policy, function<bool()> task)
        : jobId(id), name(name), priority(priority),
          status(JobStatus::PENDING), retryCount(0),
          retryPolicy(policy), task(task),
          scheduledTime(chrono::steady_clock::now()) {}

    // For deserialization (no task -- needs re-registration)
    Job(const string& id, const string& name, JobPriority priority,
        RetryPolicy policy, JobStatus status, int retryCount)
        : jobId(id), name(name), priority(priority),
          status(status), retryCount(retryCount),
          retryPolicy(policy), task(nullptr),
          scheduledTime(chrono::steady_clock::now()) {}

    string getId() const { return jobId; }
    string getName() const { return name; }
    JobPriority getPriority() const { return priority; }

    JobStatus getStatus() {
        lock_guard<mutex> lock(jobMtx);
        return status;
    }

    void setStatus(JobStatus s) {
        lock_guard<mutex> lock(jobMtx);
        status = s;
    }

    int getRetryCount() {
        lock_guard<mutex> lock(jobMtx);
        return retryCount;
    }

    void incrementRetry() {
        lock_guard<mutex> lock(jobMtx);
        retryCount++;
    }

    const RetryPolicy& getRetryPolicy() const { return retryPolicy; }

    void setLastError(const string& err) {
        lock_guard<mutex> lock(jobMtx);
        lastError = err;
    }

    string getLastError() {
        lock_guard<mutex> lock(jobMtx);
        return lastError;
    }

    bool hasTask() const { return task != nullptr; }

    void setTask(function<bool()> t) { task = t; }

    // Execute the task. Returns true on success, false on failure.
    // Called by worker threads -- job mutex NOT held during execution
    // (to avoid blocking other operations on this job).
    bool execute() {
        if (!task) return false;
        return task();
    }

    // Serialize for persistence
    string serialize() const {
        ostringstream oss;
        oss << jobId << "|" << name << "|"
            << (int)priority << "|" << (int)status << "|"
            << retryCount << "|"
            << retryPolicy.getMaxRetries() << "|"
            << retryPolicy.getBaseDelayMs();
        return oss.str();
    }

    // Deserialize
    static shared_ptr<Job> deserialize(const string& line) {
        istringstream iss(line);
        string id, name;
        int pri, stat, retry, maxRetry, baseDelay;
        char delim;
        getline(iss, id, '|');
        getline(iss, name, '|');
        iss >> pri >> delim >> stat >> delim >> retry >> delim
            >> maxRetry >> delim >> baseDelay;
        return make_shared<Job>(id, name, (JobPriority)pri,
                                 RetryPolicy(maxRetry, baseDelay),
                                 (JobStatus)stat, retry);
    }

    void display() {
        lock_guard<mutex> lock(jobMtx);
        cout << "  Job[" << jobId << "] " << name
             << " | " << priorityToString(priority)
             << " | " << statusToString(status)
             << " | retries: " << retryCount << "/" << retryPolicy.getMaxRetries();
        if (!lastError.empty()) cout << " | err: " << lastError;
        cout << endl;
    }
};

#endif
\end{lstlisting}

% ============================================================
\newpage
\section{strategies.hpp}
% ============================================================
\begin{lstlisting}[title={\texttt{strategies.hpp}}]
#ifndef STRATEGIES_HPP
#define STRATEGIES_HPP

#include "models.hpp"
#include <queue>
#include <memory>
using namespace std;

// ============================================================
// STRATEGY PATTERN: Job Scheduling Policy
// Determines the ORDER in which jobs are dispatched to workers.
// unique_ptr<SchedulingStrategy>: JobScheduler is sole owner.
// ============================================================

class SchedulingStrategy {
public:
    virtual ~SchedulingStrategy() = default;
    virtual shared_ptr<Job> selectNext(vector<shared_ptr<Job>>& pendingJobs) = 0;
    virtual string getName() const = 0;
};

// Highest priority first, then FIFO within same priority
class PriorityFirstStrategy : public SchedulingStrategy {
public:
    shared_ptr<Job> selectNext(vector<shared_ptr<Job>>& pendingJobs) override {
        if (pendingJobs.empty()) return nullptr;
        int bestIdx = 0;
        for (int i = 1; i < (int)pendingJobs.size(); i++) {
            if (pendingJobs[i]->getPriority() > pendingJobs[bestIdx]->getPriority()) {
                bestIdx = i;
            }
        }
        auto job = pendingJobs[bestIdx];
        pendingJobs.erase(pendingJobs.begin() + bestIdx);
        return job;
    }
    string getName() const override { return "PriorityFirst"; }
};

// Simple FIFO
class FIFOStrategy : public SchedulingStrategy {
public:
    shared_ptr<Job> selectNext(vector<shared_ptr<Job>>& pendingJobs) override {
        if (pendingJobs.empty()) return nullptr;
        auto job = pendingJobs.front();
        pendingJobs.erase(pendingJobs.begin());
        return job;
    }
    string getName() const override { return "FIFO"; }
};

// ============================================================
// OBSERVER PATTERN: Job Lifecycle Notifications
// shared_ptr<JobObserver>: shared between scheduler + caller.
// ============================================================

class JobObserver {
public:
    virtual ~JobObserver() = default;
    virtual void onJobCompleted(shared_ptr<Job> job) = 0;
    virtual void onJobFailed(shared_ptr<Job> job) = 0;
    virtual void onJobDeadLettered(shared_ptr<Job> job) = 0;
    virtual string getName() const = 0;
};

class LoggingObserver : public JobObserver {
public:
    void onJobCompleted(shared_ptr<Job> job) override {
        cout << "    [Log] Job " << job->getId() << " completed" << endl;
    }
    void onJobFailed(shared_ptr<Job> job) override {
        cout << "    [Log] Job " << job->getId() << " failed (retry "
             << job->getRetryCount() << "): " << job->getLastError() << endl;
    }
    void onJobDeadLettered(shared_ptr<Job> job) override {
        cout << "    [Log] Job " << job->getId() << " -> DEAD LETTER QUEUE" << endl;
    }
    string getName() const override { return "Logger"; }
};

class AlertObserver : public JobObserver {
public:
    void onJobCompleted(shared_ptr<Job> job) override {}
    void onJobFailed(shared_ptr<Job> job) override {}
    void onJobDeadLettered(shared_ptr<Job> job) override {
        cout << "    [ALERT] CRITICAL: Job " << job->getId()
             << " dead-lettered after " << job->getRetryCount()
             << " retries!" << endl;
    }
    string getName() const override { return "Alert"; }
};

#endif
\end{lstlisting}

% ============================================================
\newpage
\section{managers.hpp}
% ============================================================
\begin{lstlisting}[title={\texttt{managers.hpp}}]
#ifndef MANAGERS_HPP
#define MANAGERS_HPP

#include "strategies.hpp"
#include <unordered_map>
#include <mutex>
#include <thread>
#include <condition_variable>
#include <fstream>
#include <atomic>
using namespace std;

// ============================================================
// LOCKING: 2-level pessimistic
// Level 1: schedulerMtx (JobScheduler)
//   -> pendingJobs, allJobs, deadLetterQueue, observers, strategy
//   -> lock_guard for simple reads/writes
//   -> unique_lock + condition_variable for worker wait/notify
//
// Level 2: jobMtx (per-Job)
//   -> job status, retry count, last error
//   -> lock_guard inside Job methods (atomic status updates)
//
// WORKER POOL:
//   Workers block on condition_variable when no jobs pending.
//   submitJob signals one worker. Workers loop until shutdown.
//
// PERSISTENCE:
//   persistJobs() writes all jobs to a file.
//   loadJobs() reads them back. Tasks must be re-registered
//   (functions can't be serialized).
//
// EXPONENTIAL BACKOFF:
//   On failure: delay = baseDelay * 2^attempt, then re-queue.
//   After maxRetries: move to dead letter queue.
//
// SMART POINTERS:
//   shared_ptr<Job>         -> scheduler + workers + caller all hold refs
//   unique_ptr<Strategy>    -> scheduler sole owner
//   shared_ptr<JobObserver> -> scheduler vector + caller
// ============================================================

class JobScheduler {
    static mutex singletonMtx;
    mutex schedulerMtx;
    condition_variable cv;

    vector<shared_ptr<Job>> pendingJobs;
    unordered_map<string, shared_ptr<Job>> allJobs;
    vector<shared_ptr<Job>> deadLetterQueue;
    vector<shared_ptr<JobObserver>> observers;
    unique_ptr<SchedulingStrategy> strategy;

    vector<thread> workers;
    atomic<bool> running{false};
    string persistFile = "jobs.dat";

    JobScheduler() {
        strategy = make_unique<PriorityFirstStrategy>();
    }

    // Snapshot pattern for notifications
    void notifyCompleted(shared_ptr<Job> job) {
        vector<shared_ptr<JobObserver>> snap;
        { lock_guard<mutex> lock(schedulerMtx); snap = observers; }
        for (auto& obs : snap) obs->onJobCompleted(job);
    }

    void notifyFailed(shared_ptr<Job> job) {
        vector<shared_ptr<JobObserver>> snap;
        { lock_guard<mutex> lock(schedulerMtx); snap = observers; }
        for (auto& obs : snap) obs->onJobFailed(job);
    }

    void notifyDeadLettered(shared_ptr<Job> job) {
        vector<shared_ptr<JobObserver>> snap;
        { lock_guard<mutex> lock(schedulerMtx); snap = observers; }
        for (auto& obs : snap) obs->onJobDeadLettered(job);
    }

    // Worker thread function
    void workerLoop(int workerId) {
        while (true) {
            shared_ptr<Job> job;

            // Wait for a job (unique_lock for condition_variable)
            {
                unique_lock<mutex> lock(schedulerMtx);
                cv.wait(lock, [this] {
                    return !pendingJobs.empty() || !running;
                });

                if (!running && pendingJobs.empty()) return;

                job = strategy->selectNext(pendingJobs);
                if (!job) continue;
            }
            // schedulerMtx released -- worker executes without holding it

            job->setStatus(JobStatus::RUNNING);
            cout << "  [Worker " << workerId << "] Executing "
                 << job->getId() << " (" << job->getName() << ")" << endl;

            bool success = job->execute();

            if (success) {
                job->setStatus(JobStatus::COMPLETED);
                notifyCompleted(job);
            } else {
                job->incrementRetry();
                int retries = job->getRetryCount();
                int maxRetries = job->getRetryPolicy().getMaxRetries();
                job->setLastError("execution failed");

                if (retries >= maxRetries) {
                    // Dead letter queue
                    job->setStatus(JobStatus::DEAD_LETTERED);
                    {
                        lock_guard<mutex> lock(schedulerMtx);
                        deadLetterQueue.push_back(job);
                    }
                    notifyDeadLettered(job);
                } else {
                    // Exponential backoff: sleep then re-queue
                    int delayMs = job->getRetryPolicy().getDelayMs(retries - 1);
                    cout << "  [Worker " << workerId << "] "
                         << job->getId() << " failed, retry " << retries
                         << "/" << maxRetries << " in " << delayMs << "ms" << endl;

                    notifyFailed(job);

                    this_thread::sleep_for(chrono::milliseconds(delayMs));

                    job->setStatus(JobStatus::PENDING);
                    {
                        lock_guard<mutex> lock(schedulerMtx);
                        pendingJobs.push_back(job);
                    }
                    cv.notify_one();
                }
            }
        }
    }

public:
    static JobScheduler* getInstance() {
        lock_guard<mutex> lock(singletonMtx);
        static JobScheduler instance;
        return &instance;
    }

    void setStrategy(unique_ptr<SchedulingStrategy> s) {
        lock_guard<mutex> lock(schedulerMtx);
        strategy = move(s);
        cout << "[Scheduler] Strategy: " << strategy->getName() << endl;
    }

    void addObserver(shared_ptr<JobObserver> obs) {
        lock_guard<mutex> lock(schedulerMtx);
        observers.push_back(obs);
    }

    // Start worker pool
    void start(int numWorkers) {
        running = true;
        for (int i = 0; i < numWorkers; i++) {
            workers.emplace_back(&JobScheduler::workerLoop, this, i);
        }
        cout << "[Scheduler] Started " << numWorkers << " workers" << endl;
    }

    // Graceful shutdown: finish pending, then stop
    void shutdown() {
        {
            lock_guard<mutex> lock(schedulerMtx);
            running = false;
        }
        cv.notify_all();
        for (auto& w : workers) w.join();
        workers.clear();
        cout << "[Scheduler] Shutdown complete" << endl;
    }

    // Submit a job
    void submitJob(shared_ptr<Job> job) {
        {
            lock_guard<mutex> lock(schedulerMtx);
            allJobs[job->getId()] = job;
            pendingJobs.push_back(job);
        }
        cv.notify_one();
        cout << "[Scheduler] Submitted: " << job->getId()
             << " (" << job->getName() << ") "
             << priorityToString(job->getPriority()) << endl;
    }

    // Persist all jobs to file (survive restarts)
    void persistJobs() {
        lock_guard<mutex> lock(schedulerMtx);
        ofstream out(persistFile);
        for (auto& [id, job] : allJobs) {
            out << job->serialize() << "\n";
        }
        cout << "[Scheduler] Persisted " << allJobs.size() << " jobs to "
             << persistFile << endl;
    }

    // Load jobs from file. Tasks must be re-registered via registerTask().
    void loadJobs() {
        lock_guard<mutex> lock(schedulerMtx);
        ifstream in(persistFile);
        if (!in.is_open()) {
            cout << "[Scheduler] No persist file found" << endl;
            return;
        }
        string line;
        int count = 0;
        while (getline(in, line)) {
            if (line.empty()) continue;
            auto job = Job::deserialize(line);
            allJobs[job->getId()] = job;
            // Re-queue pending/failed jobs
            if (job->getStatus() == JobStatus::PENDING ||
                job->getStatus() == JobStatus::FAILED) {
                job->setStatus(JobStatus::PENDING);
                pendingJobs.push_back(job);
            }
            count++;
        }
        cout << "[Scheduler] Loaded " << count << " jobs from " << persistFile << endl;
    }

    // Register a task function for a loaded job
    void registerTask(const string& jobId, function<bool()> task) {
        lock_guard<mutex> lock(schedulerMtx);
        auto it = allJobs.find(jobId);
        if (it != allJobs.end()) {
            it->second->setTask(task);
        }
    }

    // Display
    void displayAllJobs() {
        lock_guard<mutex> lock(schedulerMtx);
        cout << "\n=== All Jobs ===" << endl;
        for (auto& [id, job] : allJobs) job->display();
    }

    void displayDeadLetterQueue() {
        lock_guard<mutex> lock(schedulerMtx);
        cout << "\n=== Dead Letter Queue ===" << endl;
        if (deadLetterQueue.empty()) {
            cout << "  (empty)" << endl;
            return;
        }
        for (auto& job : deadLetterQueue) job->display();
    }

    int getPendingCount() {
        lock_guard<mutex> lock(schedulerMtx);
        return pendingJobs.size();
    }

    JobScheduler(const JobScheduler&) = delete;
    JobScheduler& operator=(const JobScheduler&) = delete;
};

mutex JobScheduler::singletonMtx;

#endif
\end{lstlisting}

% ============================================================
\newpage
\section{main.cpp}
% ============================================================
\begin{lstlisting}[title={\texttt{main.cpp}}]
#include <iostream>
#include <thread>
#include <chrono>
#include <atomic>
#include "managers.hpp"
using namespace std;

int main() {
    auto* sched = JobScheduler::getInstance();

    cout << "=== Job Scheduler ===" << endl << endl;

    // Observers
    sched->addObserver(make_shared<LoggingObserver>());
    sched->addObserver(make_shared<AlertObserver>());

    // Start worker pool (3 workers)
    sched->start(3);

    // ========== 1. Submit jobs with different priorities ==========
    cout << "\n--- Submitting Jobs ---" << endl;

    // Job that always succeeds
    auto job1 = make_shared<Job>("JOB_1", "SendEmail", JobPriority::LOW,
        RetryPolicy(3, 50), []() {
            this_thread::sleep_for(chrono::milliseconds(50));
            return true;
        });

    // Job that always succeeds (high priority -- runs first)
    auto job2 = make_shared<Job>("JOB_2", "ProcessPayment", JobPriority::CRITICAL,
        RetryPolicy(3, 100), []() {
            this_thread::sleep_for(chrono::milliseconds(30));
            return true;
        });

    // Job that fails twice then succeeds (tests exponential backoff)
    static atomic<int> attempt3{0};
    auto job3 = make_shared<Job>("JOB_3", "SyncDatabase", JobPriority::HIGH,
        RetryPolicy(3, 100), []() {
            int a = ++attempt3;
            this_thread::sleep_for(chrono::milliseconds(20));
            if (a <= 2) return false;  // fail first 2 attempts
            return true;               // succeed on 3rd
        });

    // Job that always fails (tests dead letter queue)
    auto job4 = make_shared<Job>("JOB_4", "FlakyAPI", JobPriority::MEDIUM,
        RetryPolicy(2, 50), []() {
            this_thread::sleep_for(chrono::milliseconds(20));
            return false;  // always fails
        });

    // Job with high priority (tests priority ordering)
    auto job5 = make_shared<Job>("JOB_5", "GenerateReport", JobPriority::HIGH,
        RetryPolicy(1, 100), []() {
            this_thread::sleep_for(chrono::milliseconds(40));
            return true;
        });

    sched->submitJob(job1);
    sched->submitJob(job2);
    sched->submitJob(job3);
    sched->submitJob(job4);
    sched->submitJob(job5);

    // Wait for all jobs to process (including retries with backoff)
    this_thread::sleep_for(chrono::milliseconds(3000));

    // ========== 2. Persistence ==========
    cout << "\n--- Persisting Jobs ---" << endl;
    sched->persistJobs();

    // ========== 3. Display results ==========
    sched->displayAllJobs();
    sched->displayDeadLetterQueue();

    // ========== 4. Shutdown ==========
    cout << "\n--- Shutdown ---" << endl;
    sched->shutdown();

    // ========== 5. Simulate restart: load persisted jobs ==========
    // In a real system this would be a new process.
    // Here we just demonstrate the API.
    cout << "\n--- Simulating Restart (load from file) ---" << endl;
    // Note: loaded jobs need tasks re-registered since functions
    // can't be serialized. Jobs without tasks will fail execution.

    cout << "\n=== Done ===" << endl;
    return 0;
}
\end{lstlisting}

\end{document}
